% ============================================================================
% 09-formalization.tex — the Lean 4 development. Toolchain facts checked against
% lean-toolchain (v4.34.0-rc2) and lake-manifest.json (mathlib tag v4.34.0-rc2,
% rev 85e3a25e...). Updated 2026-09-20 with the toolchain migration; the prior
% pin was v4.32.0 / mathlib 3dffaf2f. The migration required no source change
% and left the axiom inventory and the native_decide disclosure unchanged --
% each re-verified at the new pin, not carried over.
% See briefs/MEMO_LEAN_4_34_MIGRATION_2026_09_20.md.
% FURTHER UPDATE 2026-09-20: open_goal_partner_eq_sqrt_s7 was CLOSED
% (Agora/Sequences/SqrtBridge.lean); the development now has ZERO sorry. The
% "blocked-on-mathlib" ruling on it was wrong -- see
% briefs/BRIDGE_GOAL_CLOSED_2026_09_20.md. Sorry-count claims updated throughout.
% Axiom inventory checked against AXIOMS.md and Agora/Axioms/*.lean.
% ============================================================================
\section{The formal development}\label{sec:formalization}

The statements labeled \textbf{(K)} in this paper are theorems of a Lean~4
development which we describe here in enough detail to be audited. We also
describe what the kernel did \emph{not} check, since a formalization paper
that reports only its successes is of no use to a referee.

\subsection{Toolchain and layout}

The development builds with Lean~4 at toolchain \texttt{v4.34.0-rc2}
\cite{deMouraUllrich2021} against \texttt{mathlib4}~\cite{mathlib4,mathlibCPP2020}
pinned at tag \texttt{v4.34.0-rc2}, commit
\texttt{85e3a25e006c35636f0e53b0e9296caca2685bc0}. (The development was
migrated to this pin from Lean~\texttt{v4.32.0} / \texttt{mathlib4}
\texttt{3dffaf2f} on 2026-09-20; no source file required any change, and every
statement in the development is unchanged, as certified by a
statement-hash lock taken before the migration and re-checked after it.) It is
organized so that each epistemic category has exactly one place to live:

\begin{center}
\begin{tabular}{@{}ll@{}}
\toprule
Directory & Contents and rule \\
\midrule
\artifact{Agora/} & The mathematics. No \artifact{sorry} anywhere. \\
\artifact{Agora/Axioms/} & The only permitted location for an
  \artifact{axiom} declaration. \\
\artifact{OpenGoals/} & The only permitted location for a \artifact{sorry};
  each is a named goal. \\
\artifact{Tests/} & Golden numeric tests against independently computed
  values. \\
\bottomrule
\end{tabular}
\end{center}

% ⚠️ T0 CORRECTION 2026-09-20 (Opus 5) — AMENDED ON T0 INSTRUCTION, same date.
% The previous wording claimed both placement rules were "enforced mechanically
% by a repository hook" and referred to a "no-sorry check". Verified false this
% session; corrected below. Record of what was wrong:
%   * "enforced mechanically by a repository hook" — the only mechanism is
%     .claude/hooks/lean_guard.sh, a Claude Code PostToolUse hook. It fires ONLY
%     when Claude edits a .lean file in a session loading .claude/settings.json.
%     A hand edit, another tool, or any `git commit` bypasses it entirely.
%   * "an axiom ... is rejected" — TRUE (the hook exits 2 on that path).
%   * "the no-sorry check excludes OpenGoals/ and nothing else" — there is NO
%     blocking no-sorry check. The hook only warns and exits 0, and the repo has
%     no .github/ directory, so there is no CI at all. (CLAUDE.md rule 3 made the
%     same "CI-enforced" claim; it is corrected on this branch.)
% The paper must not claim an enforcement mechanism stronger than the one that
% exists; that is the claim a referee is most entitled to check.
The two placement rules are enforced differently, and we state the difference
because a formalization paper that overstates its own safeguards forfeits the
audit it is asking for. The \artifact{axiom} rule is enforced mechanically: an
editor hook rejects an \artifact{axiom} declaration introduced outside
\artifact{Agora/Axioms/}. The \artifact{sorry} rule is \emph{not} mechanically
enforced --- it is maintained by review, and checked by the audit commands of
\S\ref{subsec:not-checked}, which report every occurrence and its location; at
the time of writing they report none at all. We can now be sharper about
\emph{why} the build is not the gate here, because it was tested rather than
assumed: appending a \artifact{sorry}-carrying theorem to a module and
rebuilding exits \emph{zero}, reporting success with only a
\texttt{declaration uses `sorry`} warning. What does catch it is the axiom
audit, via \artifact{sorryAx}. A reader auditing this development should
therefore run the audit and not infer a clean development from a clean build. This matters because the failure mode
being guarded against is not dishonesty but drift --- an assumption introduced
locally during a proof attempt and never surfaced.

\subsection{The results, and how they are stated}

The kernel-checked content of this paper is:

\begin{center}
\begin{tabular}{@{}lll@{}}
\toprule
Statement & Lean name & File \\
\midrule
Thm.~\ref{thm:sym2-template}, $\thetaop^3$ &
  \artifact{partner\_res3} & \artifact{Sequences/PartnerOperators.lean} \\
Thm.~\ref{thm:sym2-template}, $\thetaop^2$ &
  \artifact{partner\_res2} & \artifact{"} \\
Thm.~\ref{thm:sym2-template}, $\thetaop^1$ &
  \artifact{partner\_res1} & \artifact{"} \\
Thm.~\ref{thm:sym2-template}, $\thetaop^0$ &
  \artifact{partner\_res0} & \artifact{"} \\
Lemma~\ref{lem:magic} &
  \artifact{partner\_magic} & \artifact{"} \\
Lemma~\ref{lem:c2-collapse} &
  \artifact{cooperC2\_eq\_thetaC\_cooperC3} & \artifact{"} \\
Cor.~\ref{cor:sporadic-partners} &
  \artifact{s7\_P2\_eq}, \dots, \artifact{s18\_P0\_eq} & \artifact{"} \\
Thm.~\ref{thm:avs} &
  \artifact{P\_cleared\_eq\_zero} & \artifact{Swampland/SymSquareC3b.lean} \\
Rem.~\ref{rem:golden-sym2} &
  \artifact{symSquare\_golden\_harmonic/\_exp} & \artifact{SymSquare.lean} \\
Prop.~\ref{prop:dyadic} &
  \artifact{sqrtSeq\_dyadic} & \artifact{Sequences/FormalSqrt.lean} \\
Thm.~\ref{thm:non-integral} &
  \artifact{s10\_partner\_not\_integral}, \artifact{s18\_\dots} &
  \artifact{Sequences/PartnerIntegrality.lean} \\
Cor.~\ref{cor:s7-integral}, mechanical half &
  \artifact{partnerSeq\_s7\_recurrence} & \artifact{"} \\
Cor.~\ref{cor:s7-integral} &
  \artifact{s7\_partner\_integral} & \artifact{"} \\
Rem.~\ref{rem:pass7} &
  \artifact{s7\_partner\_integral\_pass7} & \artifact{"} \\
Thm.~\ref{thm:bridge}, step (ii) &
  \artifact{conv\_cooper\_of\_rec\_generic} &
  \artifact{Sequences/SqrtBridgeGeneric.lean} \\
Thm.~\ref{thm:bridge} &
  \artifact{partner\_eq\_sqrt}, \artifact{\_s7}, \artifact{\_s10} & \artifact{"} \\
Cor.~\ref{cor:partner-dyadic} &
  \artifact{partner\_dyadic}, \artifact{s10\_partner\_dyadic} & \artifact{"} \\
Prop.~\ref{prop:mn-lean} &
  \artifact{glue\_congruence}, \artifact{swap\_is\_fricke}, \dots &
  \artifact{Geometry/MnLattice.lean} \\
Thm.~\ref{thm:mod4} &
  \artifact{four\_dvd\_s7},
  \artifact{s7\_partner\_integral\_axiom\_free} &
  \artifact{Sequences/S7Mod4.lean} \\
Prop.~\ref{prop:embedding} &
  \artifact{B\_pullback}, \artifact{C\_pullback},
  \artifact{B\_orthogonal\_C} & \artifact{Geometry/Embedding.lean} \\
Rems.~\ref{rem:two-lattices}, \ref{rem:signature-tierA} &
  \artifact{no\_isometry\_G0N\_TN}, \artifact{sym2\_isometry\_general},
  \artifact{TN\_diagonalises} &
  \artifact{Geometry/SymSquareForms.lean}, \artifact{MnLattice.lean} \\
Prop.~\ref{prop:rho} &
  \artifact{rho\_isometry}, \artifact{rho\_mul},
  \artifact{rho\_mulVec\_period} & \artifact{Geometry/ModularAction.lean} \\
Prop.~\ref{prop:selfdual} &
  \artifact{zOf\_fricke}, \artifact{s7\_P2\_discriminant},
  \artifact{W7conj\_eq\_neg\_reflection} &
  \artifact{Geometry/SelfDual.lean}, \artifact{"} \\
\bottomrule
\end{tabular}
\end{center}

Three design decisions shape how these are stated.

\emph{$\thetaop$-form, not monic $D$-form.} Normalizing an order-3 operator to
be monic in $D$ introduces genuine rational-function coefficients, and the
rational-function type in the pinned library carries no derivative API. Rather
than improvise one, every statement above lives in \artifact{Polynomial}~$\Q$:
no division, no nonvanishing side conditions, and every proof obligation a
polynomial identity. This is why Theorem~\ref{thm:sym2-template} is stated as
four coefficient identities plus the collapse lemma rather than as an operator
equation --- the factorization in $\Q(z)[\thetaop]$ is then derived on paper,
as in \S\ref{sec:sym2}, from kernel-checked polynomial facts.

\emph{Concrete identities, not existence claims.} The relation
``$\Lthree$ is the symmetric square of $\Ltwo$'' is encoded as an equality
fixing \emph{every} coefficient of $\Lthree$ as an explicit polynomial function
of those of $\Ltwo$. It is falsifiable per candidate. The development had an
earlier generation of statements of the form ``there exists a polynomial of
degree 3'', true of any cubic and therefore content-free; those were deleted
and the episode is recorded in the repository's escalation log. The lesson
generalizes: a formalized statement is worth exactly the content of its
statement, and no amount of kernel checking supplies content the statement
lacks.

\emph{Derivatives are computed, not transcribed.} Every occurrence of a
derivative in Theorem~\ref{thm:avs} is an application of the library's
\artifact{Polynomial.derivative}, so the kernel performs each differentiation
itself. Encoding a hand-computed derivative as an opaque definition would let
a transcription error pass \artifact{ring} undetected.

\subsection{The axiom inventory, in full}\label{subsec:axioms}

The development registers exactly one load-bearing literature axiom.

\begin{center}
\begin{minipage}{0.93\textwidth}
\small
\begin{verbatim}
axiom obrien2016_theorem6_2 (c : N -> Q) (h0 : c 0 = 1) (h1 : c 1 = 2)
    (hrec : forall k : N,
      ((k : Q) + 2)^2 * c (k + 2) =
        (26 * ((k : Q) + 1)^2 + 13 * ((k : Q) + 1) + 2) * c (k + 1)
          + 3 * (3*(k : Q) + 1) * (3*(k : Q) + 2) * c k) :
    forall n, exists m : Z, c n = (m : Q)
\end{verbatim}
\end{minipage}
\end{center}

\noindent
This is Theorem~\ref{thm:obrien} verbatim, re-indexed by $n = k+1$ to avoid
truncated subtraction on $\mathbb{N}$ and with the source's boundary
computation replaced by the two initial values as hypotheses. It is used
exactly once, to conclude \artifact{s7\_partner\_integral} from
\artifact{partnerSeq\_s7\_recurrence}. It is registered in the repository's
axiom register with its provenance, and the source was read in full before
being cited --- which mattered: an earlier attribution in the program named
the wrong theorem of the same source and a mechanism (``$\eta$-quotients have
integer coefficients, hence a corollary'') that does not in fact establish
integrality, since the object is a series reversion. See
Proposition~\ref{prop:normalized-reversion} for what the correct mechanism is.

It is not provable in the pinned library: O'Brien's argument needs the
$q$-expansions of two specific weight-1 level-7 modular objects, obtained by
classical modular-forms machinery that \texttt{mathlib} at the pinned commit
does not contain. Discharging it is open question~3 of
\S\ref{sec:framework}.

For completeness: the register carries one further axiom declaration,
\artifact{pipeline\_upper\_bound}, belonging to a part of the repository
unrelated to the mathematics of this paper. It is disclosed in the register as
vacuously true --- its statement is an existential satisfied by the witness
$1$ and encodes no data --- and no file used in this paper imports it. We
mention it because an axiom inventory that reports only the axioms one wishes
to discuss is not an inventory.

\subsection{Golden validation and PASS(\texorpdfstring{$N$}{N}) discipline}

Two devices keep the development from certifying its own mistakes.

\emph{Golden tests} check a formula against an answer obtained by other means.
The symmetric-square formula~\eqref{eq:sym2-formula} is validated against two
examples whose answers are derivable by inspection
(Remark~\ref{rem:golden-sym2}); no candidate-level symmetric-square theorem is
accepted until those are green, because the formula is the single point of
failure for every such theorem. The sequence values in \artifact{Tests/} are
checked against reference values computed independently in exact integer
arithmetic outside Lean, for $n \leq 19$. Separately, the recurrence-defined
partner sequence is checked against the formal square root of the bulk
sequence, on both the integral side ($s_7$: $2, 22, 336$) and the
non-integral side ($s_{10}$: $1, \tfrac{17}{2}$) --- so the $\tfrac{17}{2}$
witnessing Theorem~\ref{thm:non-integral} is not an artifact of one encoding.

\emph{PASS($N$)} is the reporting convention for a kernel-checked finite
range. \artifact{s7\_partner\_integral\_pass7} verifies integrality for
$n\leq7$ and is cited as PASS(7), never as ``the sequence is integral'' --- a
distinction that is not pedantry here, since establishing integrality needs
every term while refuting it needs one. The convention survived the arrival of
the axiom-backed general theorem: the finite statement was kept rather than
deleted, and the summary theorem \artifact{cooper\_partner\_separation}
deliberately asserts only the two \emph{negative} halves of the separation, so
that no downstream file can cite ``only $s_7$ is integral'' as a single
uniformly-established fact.

\subsection{The three-strikes escalation discipline}
\label{subsec:three-strikes}

A lemma is not worked indefinitely. The development's house rule --- three
named strategies attempted and recorded, then escalate or name an open goal,
never grind unboundedly --- governs every stalled proof attempt, and it is
mechanical rather than a stylistic preference: an escalation path routes a
lemma from the working tier, through one tier up after three failed
attempts \emph{or} a verifier-uncheckable obstruction, to the top after a
further two failed attempts \emph{or} a genuine design ambiguity, with
de-escalation available once the top tier decomposes the problem.

The development's last open goal,
\artifact{open\_goal\_partner\_eq\_sqrt\_s7} (that the recurrence-defined
partner sequence equals the formal square root of the bulk $s_7$ series,
solution-level transport of Theorem~\ref{thm:sym2-template}), carries its
three-strikes record verbatim in \artifact{OpenGoals/PartnerIntegrality.lean}:
(1) direct strong induction on $n$, which fails because the inductive step
needs a convolution identity equivalent to the goal itself; (2) simultaneous
induction on the equality together with the convolution identity, which
reduces to a WZ-style certificate identity not derivable without an
operator-action-on-sequences bridge the pinned library lacks; (3) a Mathlib
search (\artifact{exact?}/\artifact{apply?}), which returns nothing, since no
\artifact{PowerSeries.sqrt} or holonomic-sequence API exists at the pinned
commit. All three are recorded as failures of route, not failures of effort.
The same file also records a fourth, post-ruling attempt (2026-07-26): the
underlying convolution identity was re-verified as an exact-arithmetic
identity to $n=38$ by two independent encodings --- evidence, not proof, that
the goal is true --- while a computer-algebra certificate search
(\artifact{ore\_algebra} symmetric power) produced an operator that
demonstrably does \emph{not} annihilate the convolution sequence, a mismatch
recorded in the file as unresolved rather than explained away. The goal was
carried as a named \artifact{sorry} in \artifact{OpenGoals/}, per the
placement rule of \S\ref{sec:formalization}, rather than being silently
dropped or its statement weakened.

That record is the reason we can report the sequel honestly. On 2026-09-20 the
goal was \emph{closed}, by a fifth route that none of the four recorded
attempts had considered: a convolution sum may be symmetrized in its two
indices, after which the solution-level transport is a single linear
combination against the order-2 recurrence. The ``blocked-on-library'' verdict
that attempt (3) had supported was therefore wrong --- and wrong for the second
time on this same pair of goals, the sibling
\artifact{open\_goal\_partner\_integral\_s7} having been closed earlier by a
citation the same verdict had ruled out. We draw the obvious conclusion rather
than a flattering one: the discipline's value is that it \emph{records
routes and preserves the statement}, not that its verdicts are reliable. Four
recorded failures are evidence about four routes. Because the statement had
never been weakened to accommodate them, the later proof closed the original
goal verbatim --- checked here by a statement lock against the untouched
declaration, together with a negative control confirming that the lock can
fail.

The discipline is not merely a stopping rule: the sibling goal
\artifact{open\_goal\_partner\_integral\_s7} was opened under an identical
three-strikes record (induction on the recurrence, a $p$-adic-valuation
strengthening, and a Mathlib search --- each with its own recorded failure
mode) and was later \textbf{closed} --- not by any of the three grind
strategies succeeding, but by locating the applicable literature theorem
(O'Brien's Theorem~6.2, cited as the disclosed axiom
\artifact{obrien2016\_theorem6\_2} of \S\ref{sec:formalization} above, not
re-proved) after escalation. The record of
the failed attempts was kept rather than deleted once the citation route
closed the goal, on the view that a documented dead end is itself
information for the next person who reaches the same lemma.

\subsection{What the kernel did not check}\label{subsec:not-checked}

\begin{itemize}
\item \emph{Theorem~\ref{thm:obrien}}, as discussed. Kernel-checked:
  that our sequence satisfies its hypotheses.
\item \emph{That Gorodetsky's operator is the right operator for the Cooper
  sequences.} Our $\theta$-coefficients \artifact{cooperC3}--\artifact{cooperC0}
  encode eq.~(1.7) of \cite{Gorodetsky2023}, and every symmetric-square identity
  of \S\ref{sec:sym2} is a statement about that encoding. The \emph{transcription}
  has been verified: expanding
  $\theta^3 - z(2\theta+1)(a\theta^2 + a\theta + b) + z^2(c(\theta+1)^3 + d(\theta+1))$
  and collecting powers of $\theta$ reproduces all four encoded coefficients,
  checked twice by independent routes with a negative control. What is
  \emph{not} re-derived here is that (1.7) is itself the correct operator for
  these sequences; that is Gorodetsky's result, cited. The distinction matters:
  a sign or commutator slip in the transcription would have left every identity
  in \S\ref{sec:sym2} true of an operator that is not Cooper's.
\item \emph{Everything with status \textbf{(E)} or \textbf{(N)}}:
  \S\ref{sec:irreducibility}, \S\ref{sec:lattice}, and the computational parts
  of \S\ref{sec:integrality} are Python and computer-algebra computations, not
  Lean. They are exact and reproducible (\S\ref{sec:reproducibility}) but they
  are not kernel-checked, and the paper never labels them as though they were.
\item \emph{The golden tests in} \artifact{Tests/} use
  \artifact{native\_decide} for terms up to $n=19$, whose binomial
  coefficients make kernel evaluation impractical. \artifact{native\_decide}
  trusts the compiler in addition to the kernel; both affected theorems carry
  an explicit trust annotation. No result in \S\S\ref{sec:sym2}--\ref{sec:integrality}
  depends on them.
\item \emph{No open goal remains.} The last one --- the equality, for all $n$,
  of the recurrence-defined partner sequence with the formal square root of the
  bulk sequence (\artifact{open\_goal\_partner\_eq\_sqrt\_s7}) --- was
  \textbf{closed on 2026-09-20}, and the development now contains no
  \artifact{sorry} at all. It is the solution-level transport of the operator
  identity that Theorem~\ref{thm:sym2-template} establishes at operator level.
  The proof (\artifact{Agora/Sequences/SqrtBridge.lean}) depends only on Lean's
  three standard axioms and does not use the citation axiom of
  \S\ref{subsec:axioms} below.

  We record why it had been open, because the reason is instructive and
  cautionary for this style of work. The goal had been ruled
  ``blocked-on-\texttt{mathlib}'' on the grounds that the library provides
  neither a formal-power-series square root nor a holonomic-sequence API. Both
  observations are correct and neither is relevant: the proof uses only
  \artifact{PowerSeries.mk}, \artifact{coeff\_mul}, antidiagonal reindexing,
  \artifact{isUnit\_iff\_constantCoeff}, the power-series domain instance, and
  \artifact{Finset.sum\_range\_reflect}, all of which were present in the
  library version pinned at the time the goal was declared blocked. The four
  recorded failed strategies were never refuted; they were circumvented. The
  step that had been missing is that a convolution sum may be symmetrized in
  its two indices, after which the solution-level transport --- which the
  earlier analysis had taken to be a creative-telescoping problem --- reduces to
  a single linear combination against the order-2 recurrence. A
  ``blocked-on-library'' verdict is a statement about a proof route, not about a
  goal; $n$ failed strategies are evidence about those $n$ strategies.
\end{itemize}

\subsection{An honest assessment of the formalization}

The proofs above are, as proofs, shallow: after the right statements are set
up, almost all of them are closed by \artifact{ring}. We do not claim
otherwise, and we do not think it detracts from the contribution. What is
difficult in this development is not discharging the goals but arriving at
statements that are simultaneously (i) true uniformly in the template
parameters rather than per candidate, (ii) division-free, so that they can be
stated at all in the available library, (iii) non-vacuous, in the sense that
they pin coefficients rather than assert existence, and (iv) honest about
which of them rests on a citation. The partner
of~\eqref{eq:generic-partner} being \emph{determined} by the template rather
than guessed, the collapse $\thetaop(P_2) = 2P_1$ making the identity
polynomial, and the axiom quarantine are the three places where that work
went. To our knowledge, no kernel-checked proof of a symmetric-square
structure theorem for Picard--Fuchs-type operators has previously been given
in any proof assistant.

\subsubsection*{A failure mode the gates do not catch}

One further disclosure belongs in an honest assessment, because it bears on how
much a reader should infer from the fact that this development passes its
checks. A systematic audit comparing each declaration's \emph{name} against
what its \emph{statement} actually mentions found nine instances of a single
defect: a theorem that is true, compiles, and passes the build, the
\artifact{sorry} grep, the axiom audit \emph{and} the statement lock, while
proving less than its name says --- the claim living in the identifier and the
docstring rather than in the statement. One example, since the shape is easier
to recognise than to describe: a lemma named for the primitivity of $e + Nf$ in
a hyperbolic plane had the statement $\mathrm{IsCoprime}\,(1 : \mathbb{Z})\,N$,
which is true for every $N$ in every commutative ring and mentions no vector,
no basis and no lattice.

Two things are worth saying about this. First, \textbf{none of the nine was
found by any gate}; all were found by reading statements against names, prompted
by an exchange with an independent formalization effort. A reader should weight
``passes every check'' accordingly --- the checks are necessary and they are not
sufficient, and this is a sharper statement of the vacuity caveat already made
in \S\ref{subsec:not-checked}. Second, and more usefully: \textbf{every one of
the nine occurred in the repository's legacy physics-facing modules, and none in
the arithmetic and lattice development that this paper reports}. Those modules
have since been separated into a quarantine directory that nothing in the core
imports; they are retained rather than deleted, and none of them is cited
anywhere in this paper. The results of \S\S\ref{sec:sym2}--\ref{sec:integrality}
and \S\ref{sec:framework} were subjected to the same audit and are
unaffected.

% ----------------------------------------------------------------------------
% Generated-by: Opus 5 (Stream 1 paper agent); \S9.4 (three-strikes discipline)
% added 2026-07-30 by Sonnet 5 (T0 Directive 3) | Verified-by: toolchain vs
% lean-toolchain (leanprover/lean4:v4.34.0-rc2) and lake-manifest.json (mathlib
% tag v4.34.0-rc2, rev 85e3a25e006c35636f0e53b0e9296caca2685bc0 -- re-verified
% at the new pin 2026-09-20 by Opus 5; axiom inventory, native_decide
% disclosure and sorry location all re-checked, all unchanged); axiom text transcribed from
% Agora/Axioms/OBrien2016.lean; second axiom + hook rules from AXIOMS.md;
% native_decide disclosure from Tests/CooperSequences.lean header; remaining
% sorry located by grep (OpenGoals/PartnerIntegrality.lean:201 only);
% three-strikes record transcribed verbatim from
% OpenGoals/PartnerIntegrality.lean (both goals, incl. the CLOSED one's
% original record) and the escalation-path wording cross-checked against
% CLAUDE.md rule 5 and EXECUTION_PLAN.md §"Escalation path".
% | Reviewed-by: pending T0 (Xavier)
% ----------------------------------------------------------------------------
